\subsection{INTERSECTION AND SUM OF SUBMODULES}

For every family \((\mathbf{M}_{\iota})_{\iota \in I}\) of submodules of an A-module \(\mathbf{E}\) , the intersection \(\bigcap_{\iota \in I} \mathbf{E}_{\iota}\) is a submodule of \(\mathbf{E}\) . If, for each \(\iota \in I\) , \(\phi_{\iota}\) denotes the canonical homomorphism \(\mathbf{E} \to \mathbf{E} / \mathbf{M}_{\iota}\) , \(\bigcap_{\iota \in I} \mathbf{M}_{\iota}\) is the kernel of the homomorphism \(\phi: x \mapsto (\phi_{\iota}(x))\) of \(\mathbf{E}\) into \(\prod_{\iota \in I} (\mathbf{E} / \mathbf{M}_{\iota})\) , in other words, there is an exact sequence

\[
0 \longrightarrow \bigcap_ {i \in I} M _ {i} \longrightarrow E \stackrel {\phi} {\longrightarrow} \prod_ {i \in I} (E / M _ {i}).\tag{27}
\]

The linear mapping \(\phi\) and the mapping

\[
\mathrm{E} / \left(\bigcap_ {i \in I} \mathrm{M} _ {i}\right)\rightarrow \prod_ {i \in I} (\mathrm{E} / \mathrm{M} _ {i})
\]

which is obtained by passing to the quotient, are called canonical.

In particular:

PROPOSITION 8. If a family \((\mathbf{M}_i)_{i\in I}\) of submodules of \(\mathbf{E}\) has intersection reduced to 0 \(\mathbf{E}\) is canonically isomorphic to a submodule of \(\prod_{i\in I}(\mathbf{E}/\mathbf{M}_i)\) .

Given a subset X of an A-module E, the intersection F of the submodules of E containing X is called the submodule generated by X and X is called a generating set (or generating system) of F (I, § 4, no. 3); for a family \((a_{i})_{i\in I}\) of elements of E, the submodule generated by the set of \(a_{i}\) is called the submodule generated by the family \((a_{i})\) .

An A-module is called finitely generated if it has a finite generating set.

PROPOSITION 9. The submodule generated by a family \((a_{i})_{i\in I}\) of elements of an A-module E is the set of linear combinations of the family \((a_{i})\) .

Every submodule of E which contains all the \(a_{i}\) also contains the linear combinations of the \(a_{i}\) . Conversely, formulae (1) and (2) of no. 1 prove that the set of linear combinations of the \(a_{i}\) is a submodule of E which obviously contains all the \(a_{i}\) and is therefore the smallest submodule containing them.

COROLLARY 1. Let \(u: \mathbf{E} \to \mathbf{F}\) be a linear mapping, \(\mathbf{S}\) a subset of \(\mathbf{E}\) and \(\mathbf{M}\) the submodule of \(\mathbf{E}\) generated by \(\mathbf{S}\) . Then \(u(\mathbf{M})\) is the submodule of \(\mathbf{F}\) generated by \(u(\mathbf{S})\) .

In particular, the image under u of any finitely generated submodule of E is a finitely generated submodule of F.

Remark. If \(u(x) = 0\) for all \(x \in S\) , then also \(u(x) = 0\) for all \(x \in M\) . We shall sometimes refer to this result as the ``principle of extension of linear identities'' or ``principle of extension by linearity''.

In particular, to verify that a linear mapping \(u:E\to F\) is of the form \(v\circ\phi\) , where \(v:E/M\to F\) is linear and \(\phi:E\to E/M\) is the canonical projection, it suffices to verify that \(u(S)=0\) .

COROLLARY 2. The submodule generated by the union of a family \((\mathbf{M}_{t})_{t\in I}\) of submodules of a module \(\mathbf{E}\) is identical with the set of sums \(\sum_{i\in I}x_i\) , where \((x_{i})_{i\in I}\) runs through the set of families of elements of \(\mathbf{E}\) of finite support such that \(x_{i}\in \mathbf{M}_{i}\) for all \(i\in I\) .

Clearly every linear combination of elements of \(\bigcup_{i\in I}M_i\) is of the above form: the converse is obvious.

The submodule of E generated by the union of a family \((\mathbf{M}_{\iota})_{\iota\in I}\) of submodules of E is called the sum of the family \((\mathbf{M}_{\iota})\) and is denoted by \(\sum_{\iota\in I}\mathbf{M}_{\iota}\) . If for all \(\iota\in I\) , \(h_{\iota}\) is the canonical injection \(M_{\iota}\to E\) and \(h:(x_{\iota})\mapsto\sum_{\iota}h_{\iota}(x_{\iota})\) the linear mapping of \(\bigoplus_{\iota\in I}M_{\iota}\) into E corresponding to the family \((h_{\iota})\) (no. 6, Proposition 6), \(\sum_{\iota\in I}M_{\iota}\) is the image of h; in other words, there is an exact sequence

\[
\bigoplus_ {i \in I} M _ {i} \xrightarrow {h} E \longrightarrow E / \left(\sum_ {i \in I} M _ {i}\right) \longrightarrow 0.\tag{28}
\]

COROLLARY 3. If \((\mathbf{M}_{\lambda})_{\lambda \in L}\) is a right directed family of submodules of an A-module \(\mathbf{E}\) , the sum \(\sum_{\lambda \in L} \mathbf{M}_{\lambda}\) is identical with the union \(\bigcup_{\lambda \in L} \mathbf{M}_{\lambda}\) .

\(\bigcup_{\lambda \in L} M_{\lambda} \subset \sum_{\lambda \in L} M_{\lambda}\) is always true without hypothesis; on the other hand, for every finite subfamily \((M_{\lambda})_{\lambda \in J}\) of \((M_{\lambda})_{\lambda \in L}\) , there exists by hypothesis a \(\mu \in L\) such that \(M_{\lambda} \subset M_{\mu}\) for all \(\lambda \in J\) , hence \(\sum_{\lambda \in L} M_{\lambda} \subset M_{\mu}\) and it thus follows from Corollary 2 that \(\sum_{\lambda \in L} M_{\lambda} \subset \bigcup_{\lambda \in L} M_{\lambda}\) .

COROLLARY 4. Let \(0 \to \mathbf{E} \xrightarrow{f} \mathbf{F} \xrightarrow{g} \mathbf{G} \to 0\) be an exact sequence of A-modules, S a generating system of E, T a generating system of G. If \(\mathbf{T}'\) is a subset of F such that \(g(\mathbf{T}') = \mathbf{T}, \mathbf{T}' \cup f(\mathbf{S})\) is a generating system of F.

The submodule \(\mathbf{F}'\) of \(\mathbf{F}\) generated by \(\mathbf{T}' \cup f(\mathbf{S})\) contains \(f(\mathbf{E})\) and, as \(g(\mathbf{F}')\) contains \(\mathbf{T}, g(\mathbf{F}') = \mathbf{G}\) ; hence \(\mathbf{F}' = \mathbf{F}\) .

COROLLARY 5. In an exact sequence \(0 \to \mathbf{E} \to \mathbf{F} \to \mathbf{G} \to 0\) of A-modules, if E and G are finitely generated, so is F.

PROPOSITION 10. Let M, N be two submodules of an A-module E. Then there are two exact sequences

\begin{enumerate}
\def\labelenumi{(\arabic{enumi})}
\setcounter{enumi}{28}
\tightlist
\item
\end{enumerate}

\[
0 \longrightarrow \mathrm{M} \cap \mathrm{N} \stackrel {u} {\longrightarrow} \mathrm{M} \oplus \mathrm{N} \stackrel {i - j} {\longrightarrow} \mathrm{M} + \mathrm{N} \longrightarrow 0\tag{30}
\]

\[
0 \longrightarrow \mathrm{E} / (\mathrm{M} \cap \mathrm{N}) \stackrel {v} {\longrightarrow} (\mathrm{E} / \mathrm{M}) \oplus (\mathrm{E} / \mathrm{N}) \stackrel {p - q} {\longrightarrow} \mathrm{E} / (\mathrm{M} + \mathrm{N}) \longrightarrow 0
\]

where \(i:\mathbf{M}\to \mathbf{M} + \mathbf{N},j:\mathbf{N}\rightarrow \mathbf{M} + \mathbf{N}\) are the canonical injections,

\[
p: \mathrm{E} / \mathrm{M} \rightarrow \mathrm{E} / (\mathrm{M} + \mathrm{N}) \quad and \quad q: \mathrm{E} / \mathrm{N} \rightarrow \mathrm{E} / (\mathrm{M} + \mathrm{N})
\]

the canonical surjections and where the homomorphisms u and v are defined as follow:

if \(f: \mathbf{M} \cap \mathbf{N} \to \mathbf{M} \to \mathbf{M} \oplus \mathbf{N}\) and \(g: \mathbf{M} \cap \mathbf{N} \to \mathbf{N} \to \mathbf{M} \oplus \mathbf{N}\) are the canonical injections, \(u = f + g\) , and if \(r: \mathbf{E}/(\mathbf{M} \cap \mathbf{N}) \to \mathbf{E}/\mathbf{M} \to (\mathbf{E}/\mathbf{M}) \oplus (\mathbf{E}/\mathbf{N})\) and

\[
s: \mathbf {E} / (\mathbf {M} \cap \mathbf {N}) \rightarrow \mathbf {E} / \mathbf {N} \rightarrow (\mathbf {E} / \mathbf {M}) \oplus (\mathbf {E} / \mathbf {N})
\]

are the canonical mappings, \(v = r + s\) .

We prove the exactness of (29): obviously \(i-j\) is surjective and \(u\) is injective. On the other hand, to say that \((i-j)(x,y)=0\) , where \(x\in\mathbf{M}\) and \(y\in\mathbf{N}\) , means that \(i(x)-j(y)=0\) , hence \(i(x)=j(y)=z\in\mathbf{M}\cap\mathbf{N}\) , whence by definition \(x=f(z)\) , \(y=g(z)\) , which proves that \(\operatorname{Ker}(i-j)=\operatorname{Im}u\) .

We prove the exactness of (30): clearly \(p - q\) is surjective. On the other hand, to say that \(v(t) = 0\) for \(t \in \mathbf{E} / (\mathbf{M} \cap \mathbf{N})\) means that \(r(t) = s(t) = 0\) , hence \(t\) is the class mod. ( \(\mathbf{M} \cap \mathbf{N}\) ) of an element \(z \in \mathbf{E}\) whose classes mod. \(\mathbf{M}\) and mod. \(\mathbf{N}\) are zero, which implies \(z \in \mathbf{M} \cap \mathbf{N}\) and \(t = 0\) . Finally, to say that \((p - q)(x, y) = 0\) , where \(x \in \mathbf{E} / \mathbf{M}\) , \(y \in \mathbf{E} / \mathbf{N}\) means that \(p(x) = q(y)\) , or also that there exist two elements \(z'\) , \(z''\) of \(\mathbf{E}\) whose classes mod. \(\mathbf{M}\) and mod. \(\mathbf{N}\) are respectively \(x\) and \(y\) and which are such that \(z' - z'' \in \mathbf{M} + \mathbf{N}\) . Hence there are \(t' \in \mathbf{M}\) , \(t'' \in \mathbf{N}\) such that \(z' - z'' = t' - t''\) , whence

\[
z ^ {\prime} - t ^ {\prime} = z ^ {\prime \prime} - t ^ {\prime \prime} = z.
\]

Let \(w\) be the class mod. (M ∩ N) of \(z; r(w)\) is the class mod. M of \(z\) and hence also that of \(z'\) , that is \(x\) ; similarly \(s(w) = y\) , which completes the proof that \(\operatorname{Ker}(p - q) = \operatorname{Im} v\) .

